\documentclass[11pt]{article}
\usepackage{icmlko}
\begin{document}

\title{구성인 줄 알았는데 최근성이었다 --- 채택은 서고 이름이 틀렸다}
\author{월드모델 실험실}
\date{노트 515 · 2026년 8월 4일}
\maketitle

\begin{abstract}
노트 514 가 \texttt{TRAINW}(유보 무게로 학습 행을 줄이는 쪽만 가중)를
채택하면서 그 이름을 ``도메인 구성을 채점에 맞춘다''로 붙였다. 노트 512 가
\emph{그 가중이 학습 행 평균 연도를 $+1.425$ 년 최근으로 민다}는 것을 찾고
사전 등록했다 --- \textbf{연도 분포를 고정하면 무엇이 남나.}

\texttt{TRAINW\_YEARFIX} 를 만들었다(연도별 재정규화로 연도 분포를 원래대로
돌린다. 확인: 학습 행 평균 연도 $2018.296$(끔) $\to 2019.174$(가중)
$\to \mathbf{2018.288}$(가중$+$고정)).

\begin{center}
\begin{tabular}{lrrl}
\hline
자 (기준 $=$ \texttt{TRAINW} 끔) & \texttt{TRAINW} & $+$연도 고정 & 남은 몫 \\
\hline
판 & $+0.0009$ & $-0.0005$ & --- \\
날짜 통제 판 & $-0.0002$ & $\mathbf{-0.0015}$ & --- \\
무리 안 판 & $+0.0016$ & $-0.0007$ & --- \\
\textbf{KR 만화} & $\mathbf{+0.0176}$ $20/20$ & $+0.0036$ $12/20$ & $\mathbf{20\%}$ \\
비게임 앱 & $+0.0040$ & $\mathbf{-0.0041}$ & --- \\
\hline
\end{tabular}
\end{center}

\textbf{사전 등록 반증 --- $80\%$ 가 최근성이었다.} 연도를 고정하면 다섯
자가 다 무너지고 둘(날짜 통제 판 $\cdot$ 앱)은 $0$ 밖 음수가 된다.
\end{abstract}

\begin{figure}[t]
\centering
\includegraphics[width=\textwidth]{figs/f1.pdf}
\caption{\textbf{가} --- KR 분해. \textbf{나} --- 자 다섯.
\textbf{다} --- 이름 고치기.}
\end{figure}

\section{채택은 안 되돌린다. 이름을 고친다}

노트 514 의 관문 수치는 그대로다 --- 짝 씨앗 $40$ 에서 자 다섯 중 넷이
$0$ 밖 양수였고 KR 이 $40/40$ 이었다. \textbf{그 이득은 진짜다.} 바뀐 것은
\emph{왜}다.

\begin{quote}
\textbf{옛 이름} --- ``도메인 구성을 채점 무게에 맞춘다.''

\textbf{새 이름} --- ``\textbf{학습을 $1.4$ 년 최근으로 민다.}'' 도메인
무게는 그것을 하는 \emph{수단}이었고, 구성 자체의 몫은 $20\%$ 뿐이다.
\end{quote}

노트 472 가 적었듯 \textbf{어제 만든 자로 그제 한 채택을 되돌리지
않는다} --- 다만 \textbf{설명이 틀렸으면 설명을 고친다.} 노트 $492 \cdot
493$ 을 취소했을 때와 같은 처리다(거기서는 측정 자체가 잡음이라 취소했고,
여기서는 측정이 살아 있으므로 이름만 고친다).

\section{그러면 곧바로 물음이 하나 생긴다}

수단이 목적을 대신하고 있다면 \textbf{직접 하면 더 낫지 않나} --- 도메인이
아니라 \textbf{연도로 가중}하면.

주의할 것이 있다. 노트 $446 \cdot 447$ 이 \emph{최근순 축}을 죽였다 --- 날
판 $+0.0186$ 인데 날짜 통제 판 $-0.0294$ 였다. 그런데 그것은 \textbf{예보
쪽} 최근성이었다(축으로 넣어 최근 것을 높게 찍는다). 여기는
\textbf{학습 쪽} 최근성이다(최근 행에서 더 배운다). \textbf{둘은 다른
기계다} --- 실제로 \texttt{TRAINW} 는 날짜 통제 판을 $-0.0002$ 밖에 안
건드린다.

\begin{quote}
\textbf{사전 등록} --- 학습 행을 \textbf{연도 최근성으로 직접 가중}한다
(반감기 $\tau$ 로 $0.5^{(2025-\text{연도})/\tau}$, $\tau \in \{4, 8\}$).
\textbf{현 챔피언(\texttt{TRAINW})보다 KR 이 더 오르고 날짜 통제 판이
$0$ 을 물면 \texttt{TRAINW} 를 대체한다.} 날짜 통제 판이 확실한 음수면
기각하고, 그때는 노트 446 이 겨냥한 것이 학습 쪽에도 있다는 뜻이다.
\end{quote}

\textbf{덤} --- 노트 500 이 ``새 자를 만들면 그 자에서 아무 손해도 없어야
할 대조를 같이 재라''고 했다. 이 실험이 그 예다. \texttt{TRAINW\_YEARFIX}
는 \emph{아무것도 안 사야} 정상인 대조인데, 실제로 다섯 자를 다 깎았다 ---
\textbf{연도 분포를 억지로 되돌리는 것 자체가 해롭다.} 그것이 이 측정의
가장 강한 증거다.

\section*{남은 것}

\begin{itemize}
\item \textbf{연도 직접 가중} --- 위에 사전 등록.
\item \textbf{\texttt{TRAINCAP}(만화 절반)도 같은 검사} --- 그것도 만화가
제일 옛날이라 최근성일 수 있다. 노트 509 의 KR $+0.0274$ 가 무엇인지
다시 물어야 한다.
\item \textbf{앙상블 블록 $12$}(노트 514 되돌림 조건) --- 돌고 있다.
\end{itemize}

\end{document}
